% =============================================================================
% CHAPITRE 3 - Methodologie de recherche
% Source : consigne/revue_litterature_v2_DSR.tex (lignes 1914 a 2155)
% =============================================================================

\chapter{M\'ethodologie de recherche}\label{chap:methodologie}

\section{Approche méthodologique : Design Science Research}\label{approche-muxe9thodologique--design-science-research}

Cette recherche adopte une démarche de \emph{Design Science Research}
(DSR), méthodologie établie en systèmes d\textquotesingle information par
Hevner, March, Park et Ram (2004) dans \emph{MIS Quarterly} et
opérationnalisée par le modèle de processus DSRM de Peffers, Tuunanen,
Rothenberger et Chatterjee (2007). Le DSR est la méthodologie appropriée
pour cette recherche car l\textquotesingle objectif n\textquotesingle est
pas de tester des hypothèses causales (approche positiviste) ni de
comprendre des significations subjectives (approche interprétiviste),
mais de \emph{concevoir et évaluer un artefact} résolvant un problème
organisationnel identifié (Hevner et al., 2004, p. 77).

L\textquotesingle artefact central est un framework
d\textquotesingle orchestration stigmergique pour systèmes multi-agents
LLM, conçu d\textquotesingle abord comme architecture généraliste puis
spécialisé pour la transformation de code. Suivant la taxonomie de
Hevner et al. (2004), cet artefact comprend trois niveaux :
un \emph{modèle} (l\textquotesingle architecture stigmergique définissant
types de marqueurs, règles comportementales des agents et propriétés
de coordination émergente), une \emph{instanciation} (le prototype
fonctionnel validant la faisabilité), et des \emph{principes de
conception} généralisables extraits des
cycles de conception-évaluation (Gregor, Chandra Kruse \& Seidel, 2020).

Le positionnement de la contribution suit le cadre de Gregor et Hevner
(2013) : il s\textquotesingle agit d\textquotesingle une exaptation
au sens de Gregor et Hevner,
c\textquotesingle est-à-dire l\textquotesingle application d\textquotesingle un
concept de solution mature (la stigmergie, étudiée depuis 1959) à un
domaine problématique nouveau (l\textquotesingle orchestration
d\textquotesingle agents LLM). Ce positionnement maximise la
revendication de nouveauté tout en s\textquotesingle appuyant sur des
fondations théoriques solides.

\section{Processus DSRM}\label{processus-dsrm}

Le processus de recherche suit les six activités du DSRM (Peffers et al.,
2007), avec un point d\textquotesingle entrée centré sur le problème :

\emph{Activité 1, Identification et motivation du problème.}
Les organisations investissent massivement dans la modernisation logicielle,
mais les approches actuelles échouent à l\textquotesingle échelle. Les
architectures multi-agents hiérarchiques présentent des taux
d\textquotesingle échec de 41\% à 86,7\% (Cemri et al., 2025), des surcoûts
en tokens de 4x à 220x (Gao et al., 2025), et des gains marginaux
décroissants avec les modèles frontières. Simultanément, les organisations
font face à des choix stratégiques complexes, casser des monolithes,
migrer vers des microservices, moderniser incrémentalement, qui
exigent des cadres d\textquotesingle orchestration flexibles plutôt que
des workflows rigides.

\emph{Activité 2, Définition des objectifs de la solution.}
Le framework doit satisfaire cinq objectifs mesurables correspondant aux
cinq objectifs de conception : (O1) proposer une architecture
généraliste de coordination stigmergique applicable à des tâches variées ;
(O2) démontrer l\textquotesingle émergence d\textquotesingle une
spécialisation des agents sans assignation explicite de rôles ;
(O3) comparer le framework à des baselines centralisées reproductibles
sur le benchmark TravelPlanner avec le même backbone, le même scorer
officiel et le même split d\textquotesingle évaluation ;
(O4) adapter le framework à la transformation de code et atteindre des
résultats compétitifs sur PolyMigration et SWE-bench ;
(O5) produire une trace d\textquotesingle audit complète satisfaisant
les exigences de traçabilité et de conformité réglementaire, validée
par un panel d\textquotesingle experts.

\emph{Activité 3, Conception et développement.}
Le prototype sera développé en trois itérations documentées, chaque
itération étant informée par les résultats de la précédente (conformément
au principe de recherche itérative (sixième ligne directrice de Hevner
et al., 2004). Itération 1 : test des rôles agents de base et des
marqueurs stigmergiques sur des tâches simples (TravelPlanner, benchmark
de planification généraliste). Itération 2 : raffinement de la
sémantique des marqueurs et des mécanismes de résolution de conflits sur
des migrations de complexité modérée (Amazon PolyMigration). Itération
3 : introduction de comportements émergents sophistiqués
(spécialisation, priorisation par intensité de phéromone) sur des
transformations multi-fichiers (SWE-bench). Chaque décision de
conception sera tracée à une théorie fondatrice, comme
l\textquotesingle exige la rigueur DSR.

\emph{Activité 4, Démonstration.}
Application du prototype à 2-3 scénarios concrets de transformation de
code de complexité croissante sur des projets open source, permettant
la reproductibilité.

\emph{Activité 5, Évaluation.}
L\textquotesingle évaluation suit le framework FEDS (Venable, Pries-Heje
\& Baskerville, 2016) avec une stratégie centrée sur les risques
techniques et l\textquotesingle efficacité, en trois épisodes :
(1) Benchmarks techniques (TravelPlanner pour la validation généraliste,
PolyMigration et SWE-bench pour la validation spécialisée) comparant
StigmergiAgentic à un panel de baselines organisationnelles
reproductibles partageant le même backbone et le même protocole
d\textquotesingle évaluation : baseline single-agent directe, baseline
single-agent avec raisonnement guidé, baseline \emph{self-refine},
baseline centralisée \emph{planner-executor}, et baseline
hiérarchique explicite implémentée en LangGraph Supervisor, sur les
métriques de précision officielle, coût, temps et \emph{overhead} de
coordination ; une tentative exploratoire de reproduction de
SwarmAgentic est conservée en annexe technique mais n\textquotesingle entre
pas dans le tableau principal en raison de sa non-reproductibilité dans
le protocole contrôlé ;
(2) Étude de cas sur un projet open source réaliste documentant les
comportements émergents et les patterns de coordination ;
(3) Test d\textquotesingle utilité auprès d\textquotesingle un panel de
5-8 experts (architectes logiciels, managers SI, ingénieurs IA) au sens
du framework FEDS, mesurant l\textquotesingle utilité perçue, la
faisabilité, la gouvernabilité et la valeur organisationnelle de
l\textquotesingle artefact via un questionnaire structuré à items
Likert.

\emph{Activité 6, Communication.}
Le présent mémoire constitue le vecteur principal de communication.
Un article ciblant un journal classé ABS/FNEGE est envisagé en
extension.

\section{Conformité aux sept lignes directrices de Hevner et al. (2004)}\label{conformituxe9-aux-sept-guidelines}

Le tableau suivant synthétise la conformité du design de recherche
aux sept guidelines du DSR :

\begin{longtable}[]{@{}p{0.28\linewidth}p{0.65\linewidth}@{}}
\toprule
\textbf{Ligne directrice (Hevner et al., 2004)} & \textbf{Application dans ce mémoire} \\
\midrule
\endhead
G1. Conception d\textquotesingle un artefact & Framework stigmergique (modèle +
instanciation + principes de conception) \\
G2. Pertinence du problème & Échecs documentés des MAS hiérarchiques ;
besoins industriels de modernisation flexible \\
G3. Évaluation rigoureuse & Trois épisodes FEDS : benchmarks quantitatifs,
étude de cas, test d\textquotesingle utilité d\textquotesingle experts \\
G4. Contributions à la recherche & Principes de conception généralisables ;
extension de la théorie de la coordination aux agents LLM \\
G5. Rigueur scientifique & Ancrage théorique multi-disciplinaire ;
métriques standardisées ; comparaison aux baselines \\
G6. Processus itératif de recherche & Trois itérations documentées avec
apprentissage inter-cycles \\
G7. Communication des résultats & Double audience : technique
(réplicabilité) et managériale (implications organisationnelles) \\
\bottomrule
\end{longtable}

\section{Champ d\textquotesingle étude}\label{champ-duxe9tude}

Le champ d\textquotesingle étude se situe à
l\textquotesingle intersection de trois domaines : le génie logiciel
(migration et transformation de code), les systèmes multi-agents
(orchestration et coordination), et le management des systèmes
d\textquotesingle information (gouvernance, coordination
organisationnelle et conformité). Le terrain empirique est constitué de
benchmarks reconnus (TravelPlanner, PolyMigration, SWE-bench) et de
dépôts de code open source permettant l\textquotesingle évaluation
reproductible. Le framework est conçu comme une architecture
généraliste, validée d\textquotesingle abord sur un benchmark de
planification (TravelPlanner), puis spécialisée et évaluée sur des
tâches de transformation de code.
Ce positionnement interdisciplinaire est cohérent avec la formation de
l\textquotesingle EMLV qui vise à former des managers comprenant les
enjeux technologiques et capables de gouverner
l\textquotesingle innovation.

\section{Collecte de données}\label{collecte-de-donnuxe9es}

La collecte de données s\textquotesingle inscrit dans le cadre DSRM et
porte sur trois sources complémentaires.

Source 1, Données du processus de conception.
Chaque itération de conception produira une documentation structurée :
décisions de conception et leurs justifications théoriques (théories
fondatrices), alternatives considérées et raisons de leur abandon,
résultats intermédiaires et apprentissages. Cette documentation
constitue la matière première pour l\textquotesingle extraction des
principes de conception.

Source 2, Données d\textquotesingle évaluation quantitative.
Le prototype générera des métriques sur chaque benchmark : taux de
succès par tâche et par projet, coût en tokens par transformation,
temps d\textquotesingle exécution, taux de rollback, \emph{overhead}
de coordination, couverture de tests avant/après, et traces
d\textquotesingle exécution (séquences d\textquotesingle actions des
agents, patterns de lecture/écriture des phéromones, cas
d\textquotesingle émergence ou d\textquotesingle impasse). Les
frontières de Pareto coût-précision (Kapoor et al., 2024) permettront de
situer l\textquotesingle approche par rapport aux alternatives.

Source 3, Données du test d\textquotesingle utilité d\textquotesingle experts.
Un panel de 5-8 experts (architectes logiciels, managers SI, ingénieurs
IA) évaluera l\textquotesingle artefact via un protocole structuré au
sens du framework FEDS (Venable et al., 2016) : présentation du concept
stigmergique, démonstration de l\textquotesingle artefact en exécution,
puis renseignement d\textquotesingle un questionnaire structuré à items
Likert mesurant l\textquotesingle utilité perçue, la faisabilité, la
gouvernabilité et la valeur organisationnelle. Le test d\textquotesingle utilité
ne constitue pas une collecte exploratoire qualitative : la stratégie
FEDS retenue est centrée sur l\textquotesingle évaluation de
l\textquotesingle artefact et non sur l\textquotesingle élicitation de
données primaires.

\section{Analyse des données}\label{analyse-des-donnuxe9es}

Les données quantitatives seront analysées selon plusieurs cadres
complémentaires. Les frontières de Pareto coût-précision (Kapoor et al.,
2024) permettront de situer l\textquotesingle approche stigmergique par
rapport à l\textquotesingle ensemble des alternatives existantes. La
comparaison au baseline Agentless (Xia et al., 2024) établira si la
complexité additionnelle de la coordination multi-agents est justifiée.
La taxonomie MAST (Cemri et al., 2025) servira de grille de lecture
des modes d\textquotesingle échec observés.

Les données qualitatives (traces de conception, évaluations experts)
seront analysées pour extraire des principes de conception selon le
format de Gregor et al. (2020) : « Pour [contexte/classe
d\textquotesingle utilisateurs], [caractéristique de
l\textquotesingle artefact] produira [résultat] parce que
[mécanisme/justification] ». L\textquotesingle objectif est de formuler
4-6 principes de conception généralisables au-delà du seul cas de la
transformation de code.

\section{Limites anticipées}\label{limites-anticipuxe9es}

Plusieurs limites méthodologiques sont anticipées. Le prototype sera
évalué sur un ensemble de benchmarks et de cas, mais la
généralisabilité restera conditionnée par le volume et la diversité des
scénarios effectivement testés. Le panel d\textquotesingle experts
(5-8 personnes) vise la profondeur du test d\textquotesingle utilité
FEDS, sans prétention à la représentativité statistique. Les résultats dépendront des modèles LLM
disponibles au moment du développement, dont les capacités évoluent
rapidement. Le protocole expérimental peut engendrer des coûts
d\textquotesingle API non négligeables, ce qui peut contraindre
l\textquotesingle ampleur des campagnes de test. Le contexte
mono-chercheur du mémoire limite les possibilités de triangulation par
les investigateurs. Enfin, le niveau de contribution visé est celui de
la théorie de conception émergente au sens de Gregor et Hevner (2013), c\textquotesingle est-à-dire
des principes de conception soutenus par une instanciation et des
preuves d\textquotesingle évaluation initiales, nécessitant des
recherches futures pour atteindre le statut de théorie de conception
complète. Ces limites sont reconnues et seront discutées dans
l\textquotesingle analyse des résultats.

